\section{Introducción: por qué una automotriz pasa de la conducción autónoma a la Physical AI}

Esta sección expone primero la pregunta central del episodio. Liu Xianming (刘先明) asumió la dirección del centro de conducción autónoma de XPeng (小鹏), y la atención externa se concentró en el cambio de liderazgo; sin embargo, lo que realmente vale la pena documentar de la entrevista es que él vuelve a entender la conducción autónoma dentro del marco de la Physical AI: el automóvil no es solo un medio de transporte, sino también un agente inteligente en el mundo real; la conducción autónoma no es solo un sistema de percepción, planificación y control, y mapas, sino un ciclo cerrado inteligente formado en conjunto por los datos, el modelo, la fábrica en la nube, el despliegue en el vehículo y la retroalimentación de la producción en serie.

El juicio más tajante del episodio es «Language es veneno». No quiere decir que el lenguaje carezca de valor, sino que en la cadena de entrenamiento de la conducción autónoma o de la IA física, si las señales de los sensores se traducen primero a tokens de un lenguaje intermedio y luego se decodifica la trayectoria a partir del lenguaje, se introducen supervisión humana, un cuello de botella de compresión y obstáculos para el data scaling. La ruta de Liu Xianming consiste en ir eliminando capas intermedias: eliminar la dependencia del lidar, eliminar las reglas de planificación y control, eliminar la estructura intermedia del sistema de extremo a extremo y, por último, eliminar también el language bottleneck.

\begin{importantbox}{Tesis central del episodio}
La transformación de XPeng hacia la IA no consiste en «agregarle un modelo grande al automóvil», sino en reescribir la conducción autónoma como Physical AI: con modelos más grandes, más datos, menos capas intermedias humanas y un ciclo cerrado de ingeniería más sólido, para que el automóvil aprenda, se despliegue e itere de forma continua en el mundo real.
\end{importantbox}

\begin{knowledgebox}{Nota sobre la estrategia visual}
Este video es una entrevista con encuadre fijo, sin slides, pizarrón ni demostraciones de producto. En el texto principal la portada se usa solo para identificar la fuente, y todas las imágenes del texto son diagramas conceptuales propios, que explican la evolución de la pila de software de conducción autónoma, la eliminación de Language, la fábrica de modelos en la nube, el ciclo cerrado de datos del fabricante de automóviles, la reducción de niveles jerárquicos en la organización y la misión del cambio de liderazgo.
\end{knowledgebox}

\subsection{Resumen de la sección}

El eje del EP120 es la «simplificación»: en lo técnico, eliminar capas intermedias; en lo organizativo, reducir niveles jerárquicos; en lo estratégico, situar la conducción autónoma dentro de la Physical AI. No es una simple entrevista sobre un cambio de personal, sino una metodología para la transformación de una automotriz hacia la IA.
